\begin{lemma}[Scheide replacement-chain tail is external (5)]
For every defective $\delta$, the full
$(\gamma(\delta),\delta)$ incidence component recorded in the setup meets
$V(T)$ only at its initial vertex $v^0_\delta$.
\end{lemma}
\begin{proof}
Suppose otherwise, and follow the recorded simple path from
$v^0_\delta$ to its outside end.  Let $v_1\ne v^0_\delta$ be its last
tree vertex.  It is an exit and hence lies in the base by
\noderef{ScheideStructuralExitsInBase}.  Let $f_1=v_1u_1$ be the following
$\delta$-edge and let $\gamma_j$ be the unique block color missing at
$v_1$.  The reduced bijection and elementariness give
$\gamma_j\ne\gamma(\delta)$.

First interchange $\alpha_0$ and $\gamma(\delta)$ outside $V(T)$ and
switch the terminal part of the initial
$(\alpha_0,\delta)$ path.  As in
\noderef{ScheideDefectiveNoCross}, properness and the unique uncolored edge
are preserved, $T-\gamma(\delta)$ is Tashkinov, $\delta$ is missing at
$v^0_\delta$, and $f_1$ stays $\delta$-colored; append $f_1,u_1$.
Independently interchange $\gamma(\delta)$ and $\gamma_j$ outside the
tree and switch the tail from $v_1$.  The same edge-by-edge check using
\noderef{OutsideTreeColorSwap} and
\noderef{TwoColorComponentSwapProper} preserves properness and the unique
uncolored edge, leaves $T-\gamma_j$ Tashkinov with $\delta$ missing at
$v_1$, and fixes $f^0_\delta$; append
$f^0_\delta,u^0_\delta$.

Put $Z=V(T)\setminus(U_{\gamma(\delta)}\cup U_{\gamma_j})$.  Its order is
odd, and $\delta$ is present at every vertex.  Counting incidences of the
$\delta$ matching shows that an odd number of $\delta$-edges leave $Z$.
The two displayed departures are distinct, so a third identity $f_2$ leaves
$Z$.  It is on neither switched tail and retains color $\delta$ in both
recolorings.  Its outside end cannot lie in $U_{\gamma_j}$, since
\noderef{ScheideStructuralInternalPair} would then provide an internal
$\alpha_0$-edge and allow two successive appendings beyond maximum; it
cannot lie in $U_{\gamma(\delta)}$ by
\noderef{ScheideDefectiveNoCross}.  Hence it lies outside $V(T)$.  Append
$f_2$ in both constructions.  Each resulting tree has the original
maximum number of vertices.

By \noderef{ScheideGlobalMaximumMissingCounts}, each end of the uncolored
edge misses at least two colors.  The two ends are distinct and the tree is
elementary, so at least four distinct colors are missing on the first two
base vertices.  Choose one such color $\beta$ distinct from
$\alpha_0,\gamma(\delta),\gamma_j$; it also differs from $\delta$, which
is present throughout $V(T)$.  Neither recoloring changed a $\beta$ edge
or its original missing witness.  Closedness of the first new maximum tree
forces the $\beta$-edge at the new end of $f_2$ to join it to $u_1$, since
the original closed set admits no such boundary edge.  In the second new
maximum tree $u_1$ is absent, so that same edge is appendable in a color
still missing on the tree.  The resulting larger Tashkinov tree is the
final contradiction.  Therefore the recorded replacement component has
no second tree vertex.
\end{proof}
